@DOCLANGUAGEBLOCK@

\usepackage{caption}
\usepackage{etoolbox}
\usepackage{fancyhdr}
\usepackage{titlesec}
\usepackage{colortbl}
\usepackage{float}
\usepackage{tikz}
\usetikzlibrary{calc}

% 行距 ≈ 1.5×正文（正文 11pt → 基线约 16.5pt，与 A3 一致）
\linespread{1.2}
\setlength{\parskip}{5pt}
\setlength{\parindent}{0pt}
\setlength{\emergencystretch}{4em}
\raggedbottom
% 段落拆页防孤行：禁止一页顶部只余段落末行（widow）或一页底部只有段落首行（club），
% 避免出现“术支持。”孤行字等；CJK 拆行同样受此约束。
\widowpenalty=10000
\clubpenalty=10000
\displaywidowpenalty=10000

% ===== 标题：章 14 / 节 12 / 小节~10.9，思源宋体加粗（A3 实测值）=====
% 编号由 LaTeX 自动给出（目录才可“编号独立成列”对齐 A3），格式与参考一致。
\titleformat{\section}{\fontsize{14.3}{20}\selectfont\bfseries}{\thesection}{0.35em}{}
\titleformat{\subsection}{\fontsize{12}{16}\selectfont\bfseries}{\thesubsection}{0.35em}{}
\titleformat{\subsubsection}{\fontsize{10.9}{15}\selectfont\bfseries}{\thesubsubsection}{0.35em}{}
\titleformat{\paragraph}{\fontsize{10.9}{15}\selectfont\bfseries}{\theparagraph}{0.35em}{}
% 标题前后留白对齐 A3 参考（实测视觉留白：小节/小小节 上方≈28pt、下方≈18.6pt）
\titlespacing*{\section}{0pt}{8pt}{13pt}
\titlespacing*{\subsection}{0pt}{21pt}{14pt}
\titlespacing*{\subsubsection}{0pt}{21pt}{14pt}
\titlespacing*{\paragraph}{0pt}{10pt}{6pt}
% 正文标题编号需到第 4 级（如 1.3.3.1）；目录仅列到第 3 级（tocdepth=3，下方设置）
\setcounter{secnumdepth}{4}
% 标题防孤行：标题若落在页末某阈值内，先分页再排，避免标题成为页末孤行。
% 阈值按手册分流：开发手册默认 120pt（保持 28 页；附录章开篇标题除外）；
% 产品手册经实测提高到 160pt，防止“标题＋后随大图浮体”被顶成页末孤标题。
\makeatletter
\newlength{\aifc@orphanreserve}
\setlength{\aifc@orphanreserve}{120pt}
\newcommand{\aifc@needheadspace}{%
  \ifdim\pagetotal>\dimexpr\pagegoal-\aifc@orphanreserve\relax
    \newpage
  \fi}
\pretocmd{\section}{\aifc@needheadspace}{}{\relax}
\pretocmd{\subsection}{\aifc@needheadspace}{}{\relax}
\pretocmd{\subsubsection}{\aifc@needheadspace}{}{\relax}
\makeatother
\setcounter{tocdepth}{3}
% 目录标题、“图N.”图注字样随语言代入（中文：目录/图，英文：Contents/Figure）。
\renewcommand{\contentsname}{@CONTENTSNAME@}
% ===== 目录页 =====
% 章级目录条目列宽对齐 A3：编号左缘在 x43（缩进 0em），标题自 x55 起（编号盒 1.1em）；
% 节(1.5em/2.3em)、小节(3.8em/3.2em)两级沿用 LaTeX 默认，恰为 A3 参考列(x59/84、x84/119)。
\makeatletter
% 大节(章)目录条目：前加 ~11.3pt 分组留白（对齐 A3 参考：章前 ≈27pt）；
% 编号列 0em 缩进 + 1.1em 编号盒 → 编号 x43、标题 x55。
\renewcommand*\l@section{\addvspace{11.3pt}\@dottedtocline{1}{0em}{1.1em}}
\makeatother
% “目录”标题下方留白（首行再叠加大节分组留白后总距 ≈48.7pt）
\let\aifcstdTOC\tableofcontents
\renewcommand{\tableofcontents}{%
  \begingroup
    \titlespacing*{\section}{0pt}{10pt}{23pt}%
    \aifcstdTOC
  \endgroup}

% ===== 图注：小号居中 "图N. 说明"（英文版为 "Figure N. 说明"）=====
\DeclareCaptionLabelFormat{a3}{@CAPTIONLABEL@}
\DeclareCaptionFont{a3}{\fontsize{8.5}{11}\selectfont}
\renewcommand{\figurename}{@FIGNAME@}
\captionsetup[figure]{font={a3,rm},labelfont={a3,rm},labelformat=a3,labelsep=none,skip=3pt}

% ===== 表格：booktabs 灰色横线（A3：无竖线、灰线、上下粗中间细）=====
% 灰白斑马纹底色：正文第 2、4、6…行 \rowcolor{zebra}（zebra.py 注入）
\definecolor{zebra}{gray}{0.937}
\renewcommand{\heavyrulewidth}{1.6pt}
\renewcommand{\lightrulewidth}{0.55pt}
\arrayrulecolor[gray]{0.6}
\AtBeginEnvironment{longtable}{\small}
\AtBeginEnvironment{table}{\small}
\AtBeginEnvironment{verbatim}{\footnotesize}
\AtBeginEnvironment{Shaded}{\footnotesize}

% ===== 页眉页脚（仿 A3）=====
\fancypagestyle{frontmatter}{%
  \fancyhf{}%
  \fancyfoot[L]{NEXFORM ROBOTICS}%
  \fancyfoot[R]{\thepage}%
  \renewcommand{\headrulewidth}{0.4pt}%
  \renewcommand{\footrulewidth}{0.4pt}%
}
\fancyhf{}
\fancyhead[R]{@DOCSHORT@}
\fancyfoot[L]{NEXFORM ROBOTICS}
\fancyfoot[R]{\thepage}
\renewcommand{\headrulewidth}{0.4pt}
\renewcommand{\footrulewidth}{0.4pt}
\pagestyle{fancy}
\AtBeginDocument{\hypersetup{pdftitle={@DOCSHORT@}}}
% 目录/正文超链接改为黑色（去除默认红色链接色）
\AtBeginDocument{\hypersetup{colorlinks=true,linkcolor=black,urlcolor=black,citecolor=black,filecolor=black}}

% ===== 图片控制（A3 风格：居中、统一上限）=====
% 说明：max width ≈ 0.82 版心宽；正方形/竖图按高上限 0.45 页高收缩，与大屏截图协调
\AtBeginDocument{\setkeys{Gin}{width=0.82\linewidth,height=0.42\textheight,keepaspectratio}}
% 图只放页顶/独立图页（不放页底），避免大图+正文同页时被挤出页框而裁剪
\makeatletter\AtBeginDocument{\def\fps@figure{tp}}\makeatother
% 放宽浮体占页比例，避免图片与正文叠加造成页框溢出
\AtBeginDocument{%
  \renewcommand{\topfraction}{0.92}%
  \renewcommand{\textfraction}{0.06}%
  \renewcommand{\floatpagefraction}{0.8}%
  \setcounter{topnumber}{4}\setcounter{totalnumber}{6}%
}

% ===== 正文字体 =====
% 中文版：语言块给出 xeCJK 的断行设置，字体块强制使用简体中文字形 SourceHan
% Serif SC（Noto Serif CJK SC 单地区文件），避免系统 ttc 默认取 JP 面；正文 11pt。
% Path= 由 build_a3.sh 从 LITEARM_CJK_FONT_DIR 代入。
% 英文版：不加载 xeCJK，改把等宽字体换成 DejaVu Sans Mono——手册里的
% ```text``` 目录树/链路图用到 U+2500 制表符，Latin Modern Mono 没有这些字形。
%
% 注意：占位符名不要出现在注释里，替换进去的多行文本会在注释中间断行。
@CJKFONTBLOCK@

